\section{Conclusion}\label{sec:conclusion}

This paper revisits the three-stage location--quality--price game of \citet{Economides1989} without changing its noncooperative pricing structure. At maximal locations, the correct price continuation is piecewise: sufficiently large quality gaps induce exclusion rather than the regular shared-market prices. Once that continuation is used globally, the reported symmetric quality profile $a_1=a_2=1/(3\lambda)$ fails to be a quality-subgame Nash equilibrium for $1/9<\lambda<4/27$. Moreover, the quality subgame at those locations admits the asymmetric exclusion equilibria $(1/\lambda,0)$ and $(0,1/\lambda)$ throughout the original interval $(1/9,2/9)$.

The correction therefore concerns equilibrium logic rather than a minor coefficient. A local first-order condition on the regular price branch cannot certify a global quality best response when upstream quality choices can switch the downstream pricing regime. Where the reported symmetric quality equilibrium survives and coexists with the certified exclusion equilibria, the latter yield higher total welfare in the baseline model; conditional on maximal locations they coincide with fixed-location planner optima.

The paper deliberately stops short of a complete equilibrium correspondence or a mixed-strategy reconstruction of the full three-stage game. Those are distinct questions. The result established here is sufficient for a narrower conclusion: the reported minimal quality differentiation and quality-stage uniqueness characterization at the maximal locations does not survive a global analysis of the original quality strategy set and its correct downstream price continuations.
